\documentclass{article}
\input{../../../../lib.sty}
\title{\texttt{NoiseThresholdPrivacyMap} for Gaussian noise with $L_{0,1,\infty}$ bounds}
\author{\MichaelShoemateAuthor}
\begin{document}
\maketitle
\contrib
\section{Hoare Triple}
\subsection*{Precondition}
\subsubsection*{Compiler-verified}
The types are as in the implementation signature.
\subsubsection*{Caller-verified}
None.
\subsection*{Postcondition}
The postcondition of \texttt{NoiseThresholdPrivacyMap}.
\begin{proof}
Negative scales and sensitivities are rejected. Integer differences permit flooring the $L_1$ and $L_\infty$ bounds.
Write $s=\min(l_1,l_0l_\infty)$, $b=\min(l_\infty,s)$.
If $s=0$, the nonzero maps are identical; if $s>0$ and the scale is zero, the returned bound is vacuous.
Otherwise write $s=qb+r$ with $0\le r<b$.
By the majorization argument in \texttt{l01inf\_to\_lp},
\[
\|x-x'\|_2^2\le qb^2+r^2.
\]
The Gaussian noise mechanism on the common keys therefore satisfies $\rho$-zCDP for
\[
\rho=\frac{qb^2+r^2}{2\sigma^2},
\]
rounded upward.
On each side, the unstable-key values obey the original integer sensitivity bounds.
By the postcondition of \texttt{threshold\_delta}, the probability of releasing any unstable key is at most the returned delta.
Conditioning on its complement does not affect the common-key mechanism, by independence.
This establishes the claimed approximate zCDP guarantee.
\end{proof}
\end{document}
